\input{preamble.tex}

\begin{document}

\begin{frame}[plain]
  \titlepage
\end{frame}

\section{Introduction}
\subsection{Proof of the Main Theorem}

\begin{frame}
  \frametitle{There Is No Largest Prime Number}
  \framesubtitle{The proof uses \textit{reductio ad absurdum}.}

  \begin{theorem}
    There is no largest prime number.
  \end{theorem}
  \begin{proof}
    \begin{enumerate}
      % The strange way of typesetting math is to minimize font usage
      % in order to keep the file sizes of the examples small.
      \item<1-| alert@1> Suppose $p$ were the largest prime number.
      \item<2-> Let $q$ be the product of the first $p$ numbers.
      \item<3-> Then $q$\;+\,$1$ is not divisible by any of them.
      \item<4-> Thus $q$\;+\,$1$ is also prime and greater than $p$.\qedhere
    \end{enumerate}
  \end{proof}
\end{frame}

\section{Colors}
% \sepframe{A separator frame}{for the European Commission beamer theme}

\begin{frame}
  \frametitle{European Commission colors}
  \framesubtitle{Available colors in this theme}

  \begin{enumerate}
    \item<1-> \color{ecmain}ecmain
    \item<2-> \color{ecaccentyellow}ecaccentyellow
    \item<3-> \color{ectext}ectext
    \item<4-> \color{ecbackgroundgrey}ecbackgroundgrey
    \item<5-> \color{ecbackgroundblue}ecbackgroundblue
  \end{enumerate}
\end{frame}

\end{document}
